\documentclass[preprint,12pt,authoryear]{elsarticle}

\usepackage{amsmath,amssymb,amsthm}
\usepackage{booktabs}
\usepackage{multirow}
\usepackage{xcolor}
\usepackage{graphicx}
\usepackage{tikz}
\usetikzlibrary{arrows.meta,positioning,fit,backgrounds,calc,shapes.geometric}
\usepackage{hyperref}

\newtheorem{proposition}{Proposition}
\newtheorem{remark}{Remark}
\newcommand{\placeholder}[1]{\textcolor{red}{[#1]}}
\newcommand{\norm}[1]{\left\lVert #1 \right\rVert}

\journal{Transportation Research Part E: Logistics and Transportation Review}

\begin{document}

\begin{frontmatter}

\title{Deep Reinforcement Learning with Few Exact Solves for the Carrier--Vehicle Traveling Salesman Problem}

\author{\placeholder{Author names, affiliations and corresponding author}}

\begin{abstract}
In the carrier--vehicle traveling salesman problem (CVTSP), a carrier launches and recovers a drone at continuous points while the drone visits a set of targets. Evaluating one visiting order requires solving a second-order cone program (SOCP), so training a deep reinforcement learning (DRL) policy on exact rewards for every sampled route consumes hundreds of millions of solver calls. We show that solving one order already gives a feasible solution for every other order: keep the take-off and landing points of each target and let the carrier drive between them in the new order. The cost of this solution, read from a table without a solver, is therefore an upper bound on the optimal cost of any order, and it is exact for the order that was solved. We build a training recipe around this bound: it gives a reward to every sampled route, guides solver-free local search and ranks candidate expert routes, while exact solves are reserved for a few routes per instance and for every imitation label. With the same 3.84 million exact solves, POMO trained on exact rewards is on average 3.8\% worse per instance than our policy on 529 literature benchmark instances. A POMO baseline trained with 299 million solves, 78 times as many as ours in training, remains better than our policy, by 0.21\% on the same instances and 1.13\% on synthetic test instances. At equal computing time, iterated local search and variable neighbourhood search are worse than our policy on these 529 instances, but with ten times the time they find better solutions than both learned policies. The recipe is therefore a way to obtain a learned CVTSP solver with a small solver budget, not a better solver.
\end{abstract}

\begin{keyword}
Carrier--vehicle routing \sep Drone logistics \sep Deep reinforcement learning \sep Second-order cone programming \sep Certified upper bound
\end{keyword}

\end{frontmatter}

\section{Introduction}
\label{sec:intro}

A drone with a short flight range can serve targets far beyond that range if a long-range platform carries it. A ship can launch an unmanned aerial vehicle (UAV) towards an offshore installation, keep moving, and recover the UAV further along its course. The carrier--vehicle traveling salesman problem (CVTSP) \citep{garone2011} models this operation. A carrier leaves a depot, releases a drone before each target, recovers it after the visit and returns to the depot once all targets have been served. Each release, visit and recovery is a sortie, and the carrier's travel from the landing point of one sortie to the take-off point of the next is a carrier leg. The planner chooses the visiting order and the take-off and landing point of each sortie in order to minimise mission completion time.

The problem combines two kinds of decisions: the visiting order, which is a permutation, and the take-off and landing points of the drone, which lie in the continuous plane. Once the order is fixed, the remaining problem is convex and can be written as an SOCP \citep{lobo1998}; when the order is also a decision, the problem becomes a mixed-integer SOCP. Exact methods solve it only for small instances \citep{gambella2018,erdogan2021}; the largest instances proved optimal so far have 25 targets, plus one with 35 \citep{vieira2026}. Larger instances are solved by heuristics, such as the two-stage iterative heuristic of \citet{li2025}, which compute every instance from scratch.

Deep reinforcement learning has become a practical alternative for routing problems when solutions must be produced quickly. A trained policy produces a route in a single forward pass, and methods such as POMO \citep{kwon2020} reach small gaps on the traveling salesman problem. These methods are trained by sampling millions of routes and scoring each one, which is cheap when the cost of a route is a sum of known edge lengths. In the CVTSP it is not: the carrier legs depend on where the drone takes off and lands, and these points are optimised jointly for the whole route, so scoring one route means solving an SOCP. Our reference baseline, POMO trained on an exact solver reward for every sampled route, made about $3 \times 10^{8}$ solver calls over 200 epochs. Within the range of speeds and endurance seen in training, such a policy can be reused without retraining, because these parameters enter the network as inputs. A new variant of the problem, such as one with additional constraints or a different spatial distribution of targets, may require training from scratch. The solver budget of that training then decides whether building a learned solver is worthwhile at all.

This paper asks whether most of those calls are necessary. The answer rests on one structural observation. In the fixed-order SOCP, the constraints on a sortie concern one target only, and the visiting order enters only through the carrier legs between sorties. If the sortie points from one optimal solution are kept and the order is changed, every sortie remains feasible and only the carrier legs change. The resulting value is an upper bound on the optimal cost of the new order, and it can be computed by adding up entries of a table. One exact solve thus says something about every route of the instance, including the routes that will never be solved.

\citet{erdogan2021} already use the idea of keeping sortie points fixed. In their iterated local search, a neighbouring order keeps the take-off and landing points of the current solution, so its cost can be computed without a new solve. What we have not found is the use of this idea to train a DRL policy. Self-improvement methods, such as the multi-action method of \citet{luttmann2025}, compute the exact cost of every sampled solution, which is cheap in the problems they study. In our recipe, the bound replaces this computation. It gives the reward for routes that are never solved, and it ranks candidate labels, each of which is solved exactly before it is used. Section~\ref{sec:gap} compares our method with the closest CVTSP and learning methods.

The paper makes three contributions. \placeholder{Contributions pending, under discussion with the authors. Every contribution except Proposition~1 is to be tied to an experiment that verifies it. Open: whether the metaheuristic comparison joins contribution~3 or becomes a fourth; what contribution~2 claims before the results of Groups B and D; whether contribution~3 states the pattern by instance size.}
\begin{enumerate}
\item We state and prove that one optimal fixed-order solution yields a table whose route sums bound the exact cost of every visiting order from above and match it at the solved order. The proof is short because it uses the structure of the CVTSP directly. The contribution lies in how this bound is put to work.
\item We build a training recipe around the table. Rewards for all sampled routes, local search and the ranking of candidate experts use the table without calling a solver, and exact solves are reserved for the anchor route, for one policy route and for the verification of every imitation label. The recipe uses 2.89 solver calls per training instance on average.
\item We compare the recipe with POMO trained on exact rewards for every route, at equal solver budget and at full training, and we report the results by instance size. At equal budget the recipe gives the better policy. The fully trained baseline used 78 times more solver calls in training, or about 29 times more when both runs are counted as run, and it remains better on synthetic instances from the training distribution; we quantify where and by how much.
\end{enumerate}

The practical result is a change in what a learned CVTSP solver costs to obtain. Training to epoch 200 on a single workstation took about 28 hours; the baseline run took about 215 hours on a different machine with 128 reward workers. The paper also reports where this comes at a price: the policy does not match the fully trained baseline on instances from its own training distribution, and metaheuristics with longer time budgets obtain better solutions. The contribution is a training method, and we state its practical implications as a comparison of computing budgets rather than as an operational study.

Section~\ref{sec:related} reviews related work. Section~\ref{sec:problem} defines the problem and Section~\ref{sec:method} presents the method. Section~\ref{sec:experiments} reports the experiments, Section~\ref{sec:discussion} discusses practical implications and Section~\ref{sec:conclusion} concludes.

\section{Related work}
\label{sec:related}

Three strands of work bear on this paper: solution methods for the CVTSP, DRL for routing, and learning under expensive objectives.

\subsection{The CVTSP and its extensions}

\citet{garone2011} studied the problem and proposed a two-step approach that first fixes an order and then optimises the sortie points. \citet{garone2014} studied the variant in which the vehicle visits several targets per sortie. \citet{gambella2018} gave a mixed-integer SOCP formulation and an exact algorithm that solves instances with up to 15 targets; the algorithm ranks visiting orders by a lower bound and solves fixed-order problems in that order, which is an exact-method precedent for screening orders with a cheap bound. \citet{erdogan2021} proposed a formulation with stronger lower bounds and an iterated local search whose neighbourhood moves keep the take-off and landing points of each target, so that a neighbouring order receives a feasible value without a solve. \citet{poikonen2020} studied the mothership and drone routing problem, with instances of up to 200 targets; \citet{erdogan2021} describe it as essentially the same problem as the CVTSP. \citet{fahradyan2020} extended the CVTSP to two carriers that share one vehicle, and \citet{liy2024} studied a joint routing problem with several vessels and UAVs, solved with a tabu search heuristic.

Recent work has moved the exact and heuristic frontiers. \citet{li2025} developed exact methods and a two-stage iterative heuristic, and their results supply the reference values used in this paper. \citet{vieira2026} propose a branch-and-cut algorithm that iteratively corrects the error of a linear approximation; it solves all benchmark instances with up to 25 customers to optimality, and one instance with 35 customers. \citet{ren2026} address a vessel with multiple UAVs through a mixed-integer SOCP and a parallel adaptive large neighbourhood search. CVTSP research is one branch of a wider literature on drone-assisted routing \citep{murray2015,agatz2018,otto2018,macrina2020,chung2020}. All of the CVTSP methods above solve each instance from scratch and spend their computing time per instance. A trained policy moves most of that time into training, although at test time it still solves one SOCP per distinct route it proposes.

Two observations from this literature shape our design. The first is that the fixed-order problem is the unit of exact evaluation: the exact cost of a candidate order requires one SOCP. The stronger heuristics avoid paying it for every candidate, either by keeping sortie points fixed during local search \citep{erdogan2021} or by searching on an approximate matrix \citep{li2025}. The second is that the order and the sortie points are strongly coupled, since a good order depends on where the carrier will be, and the carrier's position depends on the order. Two-stage methods handle this coupling by alternating between the two levels. Our table can be read as one such alternation step, frozen and reused across a whole search.

\subsection{Deep reinforcement learning for routing}

Most deep reinforcement learning methods for routing are developed for the travelling salesman problem (TSP) and the capacitated vehicle routing problem, in which trucks with limited load serve customers. In both problems, customers are fixed points and a route is an ordering of these points. Learning to construct routes began with pointer networks \citep{vinyals2015} and was extended to vehicle routing by \citet{nazari2018}. The attention model of \citet{kool2019} established the encoder--decoder architecture that most later work adopts. POMO \citep{kwon2020} exploits the multiple symmetric starts of a route, and Sym-NCO \citep{kim2022} and efficient active search \citep{hottung2022} add symmetry training and instance-level adaptation. Improvement-based policies \citep{wu2022} and heavy-decoder models \citep{luo2023} target better quality and larger sizes. Surveys by \citet{bengio2021} and \citet{mazyavkina2021} cover the field, and \citet{wuy2026} apply it to clustered vehicle routing in this journal. In these problems, the cost of a route is the sum of its edge lengths, and the load limit is checked by simple comparisons. A reward therefore costs one pass over the route, and no optimisation problem has to be solved.

The CVTSP combines two decisions: a visiting order, as in the TSP, and, for that order, the continuous points where the drone leaves and rejoins the carrier for each target. A visiting order alone therefore does not fix the cost of a route; the second decision requires solving an SOCP. DRL methods for truck--drone routing \citep{bogyrbayeva2023,peng2025} restrict launch and landing to network nodes, so their reward is still computed directly. Fewer learning methods address problems with continuous launch and landing points. Two studies learn policies for mothership problems, in which drones are launched from and recovered by a moving carrier. \citet{yan2026} study a truck that serves as a moving launch and recovery platform for two drones delivering to moving targets, and train a multi-agent proximal policy optimisation method. \citet{liu2026} apply proximal policy optimisation to eVTOL-and-drone routing with no-fly zones; the eVTOL chooses among a finite set of moves on a grid, and the policy is retrained for each instance. The single-drone mothership problem is the CVTSP itself \citep{erdogan2021}, so we do not claim that learning methods are absent from this class of problems. In both studies, however, the reward is computed by the simulated environment from elapsed time, distance and penalty terms, without solving an optimisation problem; the take-off and landing points result from the policy's step-by-step moves in the simulation. In the CVTSP formulation used here, the policy chooses only the visiting order, and the take-off and landing points are then optimised exactly for that order by solving an SOCP, so the cost of an order is known only after this solve. Our argument concerns the cost of training when every reward requires such a solve.

\subsection{Expensive objectives and self-improvement}

In some combinatorial problems, evaluating the objective of a single solution is itself costly, because it requires a simulation, an experiment or an optimisation solve; in the CVTSP it requires an SOCP. Learning methods for such problems either reduce the number of evaluations or replace evaluations with a cheaper estimate. \citet{kim2024} take the first approach: symmetric replay training reuses solutions that have already been evaluated through their symmetries, and thereby saves evaluations. It offers no estimate for a route that was never evaluated. The second approach learns a surrogate. \citet{sun2026} train a neural network offline to estimate drone flight times from trajectory-optimisation data and use it inside an order-first, split-second heuristic for truck and drone delivery. SurCo \citep{ferber2023} learns linear surrogate costs for nonlinear combinatorial problems. These methods learn or reuse a cost signal, and none of those we read provides a guarantee on routes outside the evaluated set. A learned surrogate can be accurate on average, but it can also be wrong in a direction that is hard to detect, and its training set has to be produced in advance. A bound derived from the solution of the instance at hand needs no training data and has a known direction of error, which makes it usable as a screening rule, although in our runs it accepts relatively few routes on its own (Section~\ref{sec:search}).

Self-improvement offers a second route to lower training cost, in which a policy trains on labels that it produces itself. \citet{luo2023} train by supervised learning on optimal solutions of small instances; \citet{luo2024} remove the need for such labels by letting the model improve its own solutions, and \citet{pirnay2024} sample several solutions from the current model and imitate the best one. \citet{luttmann2025} extend the idea to problems in which several agents, such as vehicles or machines, act together. On the problems these papers study, comparing two candidate solutions is cheap, so labels can be accepted or rejected by direct evaluation. In the CVTSP, the same test costs one SOCP per candidate. Our method keeps the self-improvement loop and uses the bound to decide which candidates deserve that cost. This changes the economics of the loop. Generating a hundred candidates is cheap, but verifying a hundred candidates is not, so the screening step becomes the part of the method that determines its budget.

\subsection{Research gap}
\label{sec:gap}

Table~\ref{tab:gap} compares the closest work on five properties, which are defined in the table note. Among the sources we read, none combines a continuous route-cost subproblem, a training signal for routes that were never solved, and a primal upper bound behind that signal. The table shows two separate groups: optimisation methods for the CVTSP and the mothership problem handle the continuous subproblem but learn no policy, whereas learning methods for routing and drone routing build routes with a policy but avoid the continuous subproblem or replace it by simulation. Two parts of our design have direct precedents in the CVTSP literature. Evaluating other orders with the sortie points of a solved solution appears in the local search of \citet{erdogan2021}, and ranking orders by a cheap bound before solving them appears in the exact method of \citet{gambella2018}. The contrast with \citet{li2025} concerns the matrix. Their approximate matrix is initialised in closed form and, after a fixed-order problem is solved, updated only on traversed arcs; the other entries keep their prior values. The matrix guides a heuristic and is regenerated for diversification. Our table is built from one primal solution in a single step, assigns a valid value to every arc and bounds the cost of every order. Given these precedents, we claim as new neither the upper bound itself (Proposition~\ref{prop:bound} in Section~\ref{sec:bound}), which follows from keeping the take-off and landing points of one solved order, nor the idea of assigning a cost to an order without solving its SOCP. Our claim concerns their use to train a DRL policy: as a reward for routes that are never solved and as a ranking of candidate imitation labels, each of which is verified exactly before it is used.

\begin{table}[t]
\centering
\caption{Closest work compared on properties relevant to this paper. Columns: (a) the route cost is the optimal value of a continuous optimisation problem (``nodes'': launch and landing restricted to network nodes; ``simulated'': cost accumulated by a simulator); (b) the signal a learning method trains on, or n/a if nothing is learned; (c) how a route whose subproblem is not solved receives a value; (d) a primal upper bound on the cost of every order, obtained from one solved solution; (e) a learned policy constructs the route.}
\label{tab:gap}
\small
\resizebox{\textwidth}{!}{\begin{tabular}{lccccc}
\toprule
Work & \begin{tabular}[c]{@{}c@{}}(a) Continuous\\ subproblem\end{tabular} & \begin{tabular}[c]{@{}c@{}}(b) Training\\ signal\end{tabular} & \begin{tabular}[c]{@{}c@{}}(c) Value for\\ unsolved routes\end{tabular} & \begin{tabular}[c]{@{}c@{}}(d) Primal bound\\ on every order\end{tabular} & \begin{tabular}[c]{@{}c@{}}(e) Learned policy\\ builds the route\end{tabular} \\
\midrule
\multicolumn{6}{l}{\textit{CVTSP and mothership optimisation}} \\
\citet{gambella2018} & Yes & n/a & Lower bound & No & No \\
\citet{poikonen2020} & Yes & n/a & Lower bound & No & No \\
\citet{erdogan2021}, ILS & Yes & n/a & Retained sortie points & Neighbours only & No \\
\citet{li2025} & Yes & n/a & Estimate & No & No \\
\citet{vieira2026} & Yes & n/a & Lower bound & No & No \\
\midrule
\multicolumn{6}{l}{\textit{Learning for routing and drone routing}} \\
\begin{tabular}[l]{@{}l@{}}\citet{kool2019};\\ \citet{kwon2020}\end{tabular} & No & Exact objective & Closed form & n/a & Yes \\
\begin{tabular}[l]{@{}l@{}}\citet{bogyrbayeva2023};\\ \citet{peng2025}\end{tabular} & No (nodes) & Exact objective & Closed form & n/a & Yes \\
\citet{liu2026} & No (simulated) & Simulated reward & Simulated & n/a & Yes \\
\midrule
\multicolumn{6}{l}{\textit{Learning under expensive evaluation}} \\
\citet{kim2024} & No & Exact, with replay & Not estimated & n/a & Yes \\
\begin{tabular}[l]{@{}l@{}}\citet{pirnay2024};\\ \citet{luo2024};\\ \citet{luttmann2025}\end{tabular} & No & Exact objective & Closed form & n/a & Yes \\
\midrule
This paper & Yes & Exact solves + bound & Retained sortie points & Yes & Yes \\
\bottomrule
\end{tabular}}
\end{table}

\section{Problem formulation}
\label{sec:problem}

Let $d \in \mathbb{R}^{2}$ be the depot and $p_{1},\dots,p_{J} \in \mathbb{R}^{2}$ the targets. The carrier moves at speed $v_{c}$ and the drone at speed $v_{d} > v_{c}$. The drone can fly for at most $E$ time units per sortie. Each target must be visited in exactly one sortie. For target $i$, the drone leaves the carrier at take-off point $s_{i}$, flies to $p_{i}$ and returns to the carrier at landing point $l_{i}$. The carrier travels from $s_{i}$ to $l_{i}$ during the sortie. The depot is both origin and destination, as in all benchmark instances we use.

\subsection{The fixed-order problem}

Let $\pi = (\pi_{1},\dots,\pi_{J})$ be a visiting order, and set $\pi_{0} = \pi_{J+1} = 0$ with $l_{0} = s_{0} = d$, so that the first leg starts and the last leg ends at the depot. The decision variables are the take-off and landing points $s_{i}, l_{i} \in \mathbb{R}^{2}$, the outbound and inbound flight times $t^{1}_{i}$ and $t^{2}_{i}$, the sortie duration $\tau_{i}$, and the carrier travel time $T_{k}$ on leg $k$. Sortie variables are indexed by target, and leg variables by position in the order. The carrier--vehicle problem for order $\pi$ is
\begin{align}
f(\pi) = \min \quad & \sum_{i=1}^{J} \tau_{i} + \sum_{k=0}^{J} T_{k} \label{eq:obj} \\
\text{s.t.} \quad & \norm{s_{i} - p_{i}} \le v_{d}\, t^{1}_{i}, && i = 1,\dots,J, \label{eq:out}\\
& \norm{l_{i} - p_{i}} \le v_{d}\, t^{2}_{i}, && i = 1,\dots,J, \label{eq:in}\\
& \norm{l_{i} - s_{i}} \le v_{c}\, \tau_{i}, && i = 1,\dots,J, \label{eq:sync}\\
& \tau_{i} = t^{1}_{i} + t^{2}_{i}, \quad \tau_{i} \le E, && i = 1,\dots,J, \label{eq:end}\\
& \norm{s_{\pi_{k+1}} - l_{\pi_{k}}} \le v_{c}\, T_{k}, && k = 0,\dots,J, \label{eq:leg}\\
& t^{1}_{i},\, t^{2}_{i},\, \tau_{i},\, T_{k} \ge 0 . \label{eq:nonneg}
\end{align}
Every constraint is a second-order cone or a linear constraint, so the problem is an SOCP \citep{lobo1998}. Constraints \eqref{eq:out} and \eqref{eq:in} bound the flight times from below by the flight distances, so the drone may fly more slowly or wait in the air; $\tau_{i}$ is therefore the time between take-off and landing, and \eqref{eq:sync} requires the carrier to reach $l_{i}$ from $s_{i}$ within that time. At an optimal solution, $\tau_{i}$ equals the larger of the drone's and the carrier's travel times in the sortie, as in \citet{erdogan2021}. The CVTSP asks for $\min_{\pi} f(\pi)$. Mission completion time is the objective in \eqref{eq:obj}, because the carrier is always either on a leg or within a sortie. Table~\ref{tab:notation} lists the main notation.

\begin{table}[t]
\centering
\caption{Main notation.}
\label{tab:notation}
\small
\begin{tabular}{ll}
\toprule
Symbol & Meaning \\
\midrule
$J$ & Number of targets \\
$d$, $p_{i}$ & Depot and target locations \\
$v_{c}$, $v_{d}$ & Carrier and drone speeds \\
$E$ & Maximum sortie duration \\
$\pi$ & Visiting order, with $\pi_{0} = \pi_{J+1} = 0$ denoting the depot \\
$s_{i}$, $l_{i}$ & Take-off and landing points of the sortie for target $i$ \\
$t^{1}_{i}$, $t^{2}_{i}$, $\tau_{i}$ & Outbound time, inbound time and duration of a sortie \\
$T_{k}$ & Carrier travel time on leg $k$ \\
$f(\pi)$ & Optimal value of the fixed-order problem for $\pi$ \\
$q$ & Anchor route, solved exactly \\
$C^{q}$, $F^{q}$ & Directed table and total sortie time from the solution for $q$ \\
$U_{q}(\pi)$ & Certified upper bound on $f(\pi)$ \\
\bottomrule
\end{tabular}
\end{table}

\subsection{Cost structure}

Two features of this formulation drive the method. First, constraints \eqref{eq:out}--\eqref{eq:end} involve the variables of a single target and make no reference to $\pi$. Second, the order appears only in \eqref{eq:leg}, which links the landing point of one target to the take-off point of the next. Any solution that fixes all sortie variables therefore induces a closed-form cost for every order, since the leg times are the only remaining quantities. Evaluating a fixed-order problem with $J$ targets takes tens of milliseconds at moderate sizes in our implementation, about 0.03\,s per call at $J \approx 60$, while evaluating a table sum takes microseconds. The number of orders that can be scored in one solve time differs by several orders of magnitude, which is the leverage the method relies on. This is the structure that the next section exploits. The policy in Section~\ref{sec:method} selects $\pi$, and the SOCP determines the sortie points, so the method follows the natural decomposition into a discrete upper level and a continuous lower level. Only the upper level is learned.

\section{Methodology}
\label{sec:method}

\subsection{Overview}

Figure~\ref{fig:method} shows one training step, and Table~\ref{alg:train} lists it step by step for a batch of instances. A frozen reference policy samples 200 candidate routes per instance. A small edge predictor selects one of them as the anchor, and the anchor is solved exactly. The solution yields a table $C$ that scores every candidate. Local search on the table produces an expert route, which is verified exactly before it can serve as an imitation label. The learning policy then receives two updates: a reinforcement step on rewards drawn mainly from the table, and an imitation step on the verified experts. The number of exact solves per instance is independent of the number of POMO routes, which is what keeps the solver budget small. In POMO with an exact reward, a batch of $B$ instances with $P$ routes costs $BP$ solves. Here it costs at most $3B$ (one anchor, one policy route and at most one expert check per instance), plus at most $2B$ in refresh rounds for instances with $J \ge 60$, regardless of $P$; the training run averaged 2.89 solves per instance.

\input{fig_method}

\begin{table}[t]
\caption{One training step for a batch of instances. Exact solves are marked.}
\label{alg:train}
\small
\hrule\smallskip
\noindent For each instance in the batch:
\begin{enumerate}
\setlength{\itemsep}{0pt}
\item Sample 200 routes from the frozen reference policy.
\item Rank them with the edge predictor and choose the anchor $q$.
\item \textbf{Exact solve:} solve the SOCP for $q$ and build the table $C^{q}$.
\item Score all sampled routes with $U_{q}$ and select 30 starts.
\item Run 2-opt on $C^{q}$ from each start, and LKH from a share of the starts when $J \ge 60$.
\item \textbf{Exact solve:} if the best route has $U_{q}$ below $f(q)$, solve it. Otherwise, with the near-miss probability, solve the closest route.
\item \textbf{Exact solves:} for $J \ge 60$ after epoch 60, run up to two refresh rounds of the search, with one solve per round.
\end{enumerate}
\noindent For the batch as a whole:
\begin{enumerate}
\setlength{\itemsep}{0pt}
\setcounter{enumi}{7}
\item Decode the POMO routes of the learning policy and score them with the tables.
\item \textbf{Exact solve:} solve the best-scoring route of each instance.
\item Update the policy by REINFORCE on the hybrid reward.
\item Update the policy by cross-entropy on the verified expert routes.
\item Update the edge predictor on the tables.
\end{enumerate}
\hrule
\end{table}

\subsection{A certified upper bound on all visiting orders}
\label{sec:bound}

Let $q$ be a visiting order and let $(s^{*},l^{*},\tau^{*},T^{*})$ be an optimal solution of the fixed-order problem for $q$, with each sortie variable indexed by target. Write $F^{q} = \sum_{i} \tau^{*}_{i}$ and define a directed table on the nodes $\{0,1,\dots,J\}$, with node $0$ the depot:
\begin{equation}
C^{q}_{0j} = \frac{\norm{s^{*}_{j} - d}}{v_{c}}, \qquad
C^{q}_{i0} = \frac{\norm{d - l^{*}_{i}}}{v_{c}}, \qquad
C^{q}_{ij} = \frac{\norm{s^{*}_{j} - l^{*}_{i}}}{v_{c}} .
\end{equation}
For any order $\pi$, define the route value
\begin{equation}
U_{q}(\pi) = F^{q} + \sum_{k=0}^{J} C^{q}_{\pi_{k}\pi_{k+1}}, \qquad \pi_{0} = \pi_{J+1} = 0 .
\end{equation}

\begin{proposition}
\label{prop:bound}
For every visiting order $\pi$, $f(\pi) \le U_{q}(\pi)$. Moreover, $f(q) = U_{q}(q)$.
\end{proposition}

\begin{proof}
Fix $\pi$. Keep $s^{*}_{i}$, $l^{*}_{i}$, $t^{1*}_{i}$, $t^{2*}_{i}$ and $\tau^{*}_{i}$ for each target and set $T_{k} = C^{q}_{\pi_{k}\pi_{k+1}}$. Constraints \eqref{eq:out}--\eqref{eq:end} involve only these sortie variables and do not depend on the order, so they hold because they hold for $q$. Constraint \eqref{eq:leg} holds with equality by construction, and non-negativity is immediate. The point is thus feasible for the fixed-order problem of $\pi$, with objective value $U_{q}(\pi)$. Since $f(\pi)$ is the minimum over feasible points, $f(\pi) \le U_{q}(\pi)$. For $\pi = q$, each $T^{*}_{k}$ is at its lower bound at an optimal solution because the objective strictly increases in $T_{k}$, which gives $U_{q}(q) = f(q)$.
\end{proof}

The proof does not use the fact that the route ends at the depot. If the destination differs from the origin, as the general problem allows \citep{erdogan2021}, the proposition holds with the arcs into node~$0$ measured to the destination. Three conditions apply in practice. The anchor must be solved to optimality; a merely feasible solution would still give a valid upper bound, but the equality at $q$ could fail. In the implementation, every anchor solution passes an audit of all constraints with an absolute tolerance of $10^{-5}$ before its table is used. A small violation in one sortie can only be repaired by lengthening times, so the bound is valid up to solver feasibility tolerances accumulated over the $J$ sorties, not up to $10^{-5}$ itself. The training code also checks the inequality, with a tolerance of $2 \times 10^{-5}$, whenever an expert, near-miss or refresh route is solved exactly, and it stops on a violation, which would indicate a numerical or implementation fault. The final run completed without a stop.

\begin{remark}
\label{rem:combine}
When several orders of an instance have been solved, each solution gives its own table, and the minimum of their route values, $\min_{q} U_{q}(\pi)$, is again an upper bound on $f(\pi)$. The training step takes this minimum over the anchor and expert tables (Section~\ref{sec:learning}). An entry-wise minimum of the tables is not valid: it splices together sortie points from different solutions, so that the take-off and landing points of one target could differ between its incoming and outgoing arcs.
\end{remark}

When the depot is both origin and destination, an order and its reverse have the same optimal cost, $f(\overleftarrow{\pi}) = f(\pi)$, where $\overleftarrow{\pi}$ visits the targets of $\pi$ in reverse order: exchanging the take-off and landing points of every sortie, and its outbound and inbound times, turns a solution for one into a solution for the other. The table does not share this symmetry. Its entry $C^{q}_{ij}$ runs from the landing point of $i$ to the take-off point of $j$, so reversing a route changes its internal arcs, and $U_{q}(\overleftarrow{q})$ generally exceeds $f(q)$. Local search must therefore account for the reversal of internal arcs. A 2-opt move reverses a section of the route. In a symmetric problem only the two arcs at the ends of the section change, but here every arc inside the section is traversed in the opposite direction and changes its cost. Prefix sums of the forward and reverse arc costs along the current route give this change in constant time, so each 2-opt move is still scored in constant time, without any solver call.

\subsection{Policy, anchor selection and solver-free expert search}
\label{sec:search}

The learning policy is an attention encoder--decoder of the POMO type \citep{kwon2020}, with 4 encoder layers, embedding dimension 256, 16 attention heads and feed-forward dimension 512. Each node has six features: its coordinates relative to the depot, divided by the largest depot--target distance of the instance; its distance to the depot on the same scale; a depot indicator; the speed ratio $v_{c}/v_{d}$; and the drone reach $v_{d}E$ on the same scale. The speed and endurance features let one network serve instances with different operating parameters. The policy decodes one route per POMO start, and in training every target serves as a start, so $P = J$.

Two policies interact during training. The learning policy is updated at every step. A frozen reference policy generates search material and is replaced only when the learning policy obtains a strictly better exact mean on a fixed validation set of 1,000 instances. This gate ensures that the expert search starts from routes of validated quality.

An edge predictor maps the detached encoder output to a non-negative directed table of size $(J+1) \times (J+1)$. It is trained with a Huber loss towards tables built from exact solutions, and its loss does not reach the policy. Its only role is to rank the 200 reference routes and choose the anchor, so a poor prediction costs search quality and cannot affect validity.

The anchor table scores all sampled routes. We select 30 starts: the 20 routes with the lowest bound, and 10 further routes chosen for diversity from a pool of the 60 next-best routes within a relative gap of 5\%. From each start, directed 2-opt runs on the table for at most 30 moves and accepts only strict decreases of $U_{q}$. Local search on the table is complementary to LKH \citep{lin1973,helsgaun2000}: on large instances, a share of the starts is also passed to LKH as an asymmetric TSP on $C^{q}$. In every case the search minimises a surrogate that can only overestimate the true cost. The LKH time budget for all participating starts is deterministic, $1.5 \times 10^{-4} J^{2}$ seconds scaled by the LKH share, so that the search cost of a batch can be predicted in advance.

Proposition~\ref{prop:bound} gives the screening rule. If a route satisfies $U_{q}(\pi) < f(q)$, then $f(\pi) < f(q)$ holds without any solve, so the route is a certified improvement. Its exact cost is then computed once, and the route becomes the expert. If no route clears the threshold, the anchor remains the expert. Because the bound is conservative, some truly better routes fail this test. With a small probability, the near-miss rule therefore solves the route with the smallest relative gap above the threshold and accepts it if it is exactly better. The probability adapts to training progress: for small instances it moves between 1\% and 5\% as the imitation accuracy and the share of policy routes that beat the expert rise. Each instance has a budget of one such verification, and the solve is charged whether or not the route is accepted. If the exact cost of a near-miss route exceeds its bound, the run stops, because that would indicate a numerical or implementation fault.

In practice the first-round test accepts relatively few experts once training has progressed, and most accepted experts come from near-miss checks and refresh rounds (Section~\ref{sec:recipe}). Through epoch 133, the first-round test led to 128,154 exact verifications, near-miss checks accepted 159,866 routes and refresh rounds accepted 162,888. At epoch 133, 659 of 10,000 instances passed the first-round test, whereas 3,788 received an improved expert. Near-miss routes checked at that epoch had bounds 14.6\% above the exact cost of the anchor on average, and 40.8\% of them were exactly better than the anchor. The bound therefore works less as a certificate than as a ranking that decides which few routes receive an exact solve. Its guarantee matters for validity: every accepted label is exact, whichever rule proposed it.

\subsection{Learning from few exact solves}
\label{sec:learning}

Each instance yields three kinds of cost information for the POMO routes $\pi^{1},\dots,\pi^{P}$ of the learning policy. The anchor table and the expert table bound every route, and the code takes the route-wise minimum of the two. The route with the lowest bound is then solved exactly. Writing $c^{p}$ for the exact cost of that route and for the bound of every other route, the hybrid advantage is
\begin{equation}
a^{p} = \frac{\bar{c} - c^{p}}{\max\{\sigma(c),\, 10^{-6}\}},
\end{equation}
where $\bar{c}$ and $\sigma(c)$ are the mean and standard deviation over the POMO routes of the instance. Standardising within each instance prevents hard instances from dominating the gradient. A second term compares the exactly solved policy route $\pi^{\mathrm{sel}}$ with the verified expert $\pi^{E}$,
\begin{equation}
a^{\mathrm{ex}} = \operatorname{clip}\!\left(\frac{f(\pi^{E}) - f(\pi^{\mathrm{sel}})}{|f(\pi^{E})|},\, -1,\, 1\right).
\end{equation}
The reinforcement loss combines the two terms using REINFORCE \citep{williams1992}:
\begin{equation}
\mathcal{L}_{\mathrm{RL}} = -\frac{1}{BP}\sum_{b,p} a^{p}_{b} \log p_{\theta}(\pi^{p}_{b}) - \frac{w}{B} \sum_{b} a^{\mathrm{ex}}_{b} \log p_{\theta}(\pi^{\mathrm{sel}}_{b}),
\end{equation}
where $B$ is the batch size. The weight $w$ rises linearly from 0 to 1 as the imitation accuracy of the previous epoch moves from 0.50 to 0.85, so the comparison with the expert begins only when the policy can produce competitive routes.

The hybrid advantage $a^{p}$ is biased, not only noisy. The solved route carries its exact cost while every other route carries a bound that exceeds its exact cost by a wide margin, so the solved route tends to receive the largest advantage by construction. In a diagnostic on 200 instances with two checkpoints of an earlier run of the method (epochs 100 and 365, not E133), the cosine between this advantage vector and the one obtained with all routes solved exactly was 0.61 and 0.67, falling to 0.53 and 0.54 at $J = 100$, and the route that was truly best received a negative advantage on 8\% and 12\% of the instances. Because the bound is tightest near the anchor, the bias can also favour routes that resemble the anchor; Section~\ref{sec:limits} returns to this point.

A separate optimiser step performs teacher-forced cross-entropy on the expert sequences. Each batch mixes current experts with a replay buffer of 50,000 experts, which supplies 30\% of the batch. Only verified experts enter this update, so an error in the bound can waste a solve but can never become a label. Imitation runs for five inner epochs at each update on batches of 1,000 sequences, split into micro-batches of 128. The two updates use a common network and take separate optimiser steps.

\subsection{Progressive search}
\label{sec:recipe}

Search effort is scheduled by instance size and training stage, because a strong expert helps large instances and hurts small ones. Earlier runs showed that small instances reach their best validation value early and then deteriorate, a pattern that coincided with falling route diversity in the POMO output. Instances with $J < 60$ therefore keep the plain 2-opt search throughout. For $J \ge 60$, a ramp $r = \operatorname{clip}((e - 40)/60,\, 0,\, 1)$ at epoch $e$ raises the fraction of LKH starts to $r/2$ and raises the near-miss probability towards 1. After epoch 60, large instances also receive two refresh rounds: the most recently solved route supplies a new table, which is tight around that route, and the search continues on it. Each round solves at most one additional route.

The quantity of imitation labels is normalised at the same time. An improved expert is repeated in the imitation batch so that the expected number of samples per instance stays near 1.44, with the repeat weight clipped to $[1, 5]$. Without this step, the imitation pressure would grow as the search improves.

The values of these settings came from earlier development runs. We compared the expert operators on a frozen policy and found that swap, 2-opt and LKH improved different instances. In 20-epoch training runs, however, LKH alone was worse than swap combined with 2-opt from epoch 15 onwards, by 0.076 in mean validation cost at epoch 20, and the final recipe uses 2-opt and LKH without swap. Those runs also showed that the benefit of a stronger expert concentrates at $J \ge 60$, which set the size threshold. The ramp over epochs 40 to 100 was chosen so that the policy first learns from the cheapest search and receives the stronger expert only once it has stabilised. We did not run a formal search over these values, and Section~\ref{sec:discussion} returns to this point.

\section{Computational experiments}
\label{sec:experiments}

\subsection{Setup}

\paragraph{Baseline.} G200 is POMO trained for 200 epochs with an exact Gurobi reward on every sampled route, on 100,000 random instances per epoch drawn from the same instance generator as the proposed method. Its logged solver calls show that it decoded $J/4$ POMO starts per training instance during these 200 epochs, whereas the proposed method trains with $J$ starts; both models use the same start rule at test time. G200 uses a smaller network (6 encoder layers, embedding dimension 128, 8 heads, 1.27 million parameters, against 2.44 million for E133), a batch size of 512, and Adam with a constant learning rate of $5 \times 10^{-5}$ and weight decay $10^{-6}$. The two models therefore differ in architecture, training length and start count as well as in training signal, and the comparison in this paper is between complete training recipes. G200 was not retrained for this paper. A retraining under its recipe on our hardware is reported in Sections~\ref{sec:efficiency} and~\ref{sec:ablation}.

\paragraph{Proposed model.} Training ran for 200 epochs of 10,000 instances each, with a batch size of 64, Adam with learning rate $10^{-4}$ and weight decay $10^{-6}$, and a cosine schedule with a 500-epoch horizon that was stopped at epoch 200. We report the checkpoint at epoch 133 (E133), selected by mean exact cost on a validation set of 2,000 instances, with 400 at each of $J \in \{20,40,60,80,100\}$. The selection used no test data. One fixed-order SOCP is solved by Gurobi with one thread per call. Experiments ran on a workstation with an Intel i9-13900KF processor (32 hardware threads), 31\,GB of memory and an NVIDIA RTX 4090 GPU. Over the selection range, validation means were close: E133 reached 17.703 and E150 reached 17.705.

\paragraph{Data.} Training instances have $J$ drawn uniformly from $\{20,\dots,100\}$ and targets drawn uniformly from the unit square. Speeds and endurance are sampled uniformly between the smallest and largest values in the literature benchmark, with $v_{c} = 0.2$, $v_{d} \in [0.4, 1]$ and $E \in [0.05, 1]$ in normalised units, so the speed ratio $v_{c}/v_{d}$ lies between 0.2 and 0.5. We evaluate on three sets. Test500 contains 100 new instances at each of $J \in \{20,40,60,80,100\}$. The generalisation set contains 100 new instances at each of $J \in \{150,200,300,500\}$. Real529 is the literature benchmark of \citet{gambella2018}, \citet{erdogan2021} and \citet{poikonen2020}, as compiled by \citet{li2025}, with 72, 132 and 325 instances. Its sizes include $J$ from 10 to 19 and $J = 200$, which lie outside the training range. It contains only two parameter settings, both at the edges of the training range: 325 instances with speed ratio 0.5 and endurance 20, and 204 with speed ratio 0.2 and endurance 1, in the units of the benchmark files. Test500 serves as the primary evidence because it matches the training distribution and allows a clean comparison at each size. We inspected test500 once, for the checkpoint E150 chosen by an earlier and coarser selection, before the refined selection that chose E133; the selection itself relied on the validation set only. Real529 is not a held-out set, because earlier checkpoints of this project were evaluated on it during method development.

\paragraph{Test protocol.} Both models decode greedily with POMO, with $\min(J,100)$ starts on every set: $J$ starts up to $J = 100$ and 100 starts on larger instances, including the generalisation set and the $J = 200$ instances of Real529. Under eight-view augmentation, each geometric view is decoded separately. Every distinct route is solved exactly and the best value is reported. Costs are therefore exact in all tables.

\paragraph{Metrics.} For model $A$ and baseline $B$ on instance $n$, the relative gap is $100\,(c^{A}_{n} - c^{B}_{n})/c^{B}_{n}$, so positive values mean that $A$ is worse. We report the mean over instances with a 95\% normal-approximation interval, together with the counts of wins, ties and losses at a tolerance of $10^{-3}$. We prefer the per-instance mean because a ratio of mean costs is dominated by the largest instances. On Real529, for example, the ratio of means is +0.12\%, whereas the per-instance mean is +0.21\%.

\subsection{Results on synthetic instances}
\label{sec:synthetic}

Table~\ref{tab:test500} reports test500. Without augmentation, E133 is behind G200 at every size. The gap is +1.08\% at $J = 20$, rises to about +3.4\% at $J = 60$ and $J = 80$, and is +3.05\% at $J = 100$; the mean over all sizes is +2.60\%. G200 wins 407 of 500 instances. This is the largest gap in our evaluation, and it marks the main quality cost of the proposed training scheme. The gap does not grow steadily with size. It is smallest at $J = 20$ and stabilises between $J = 60$ and $J = 100$, which is the range in which the progressive search is active. The mean costs at $J = 20$ differ by about 0.12 units, and ties are common there, whereas at $J \ge 60$ nearly all instances are lost by a modest margin. The picture is one of a consistent shortfall and not of occasional bad instances: E133 wins only 10 to 14 instances out of 100 at the three largest sizes.

\begin{table}[t]
\centering
\caption{Test500 without augmentation. Gap is E133 relative to G200 (positive means E133 is worse).}
\label{tab:test500}
\small
\begin{tabular}{rrrrrr}
\toprule
$J$ & E133 & G200 & Gap (\%) & 95\% interval & Win/Tie/Loss \\
\midrule
20 & 12.594 & 12.472 & +1.08 & [+0.72, +1.43] & 15/22/63 \\
40 & 15.923 & 15.633 & +2.00 & [+1.55, +2.45] & 18/1/81 \\
60 & 18.502 & 17.946 & +3.39 & [+2.77, +4.01] & 13/0/87 \\
80 & 20.079 & 19.453 & +3.47 & [+2.98, +3.96] & 10/0/90 \\
100 & 22.040 & 21.440 & +3.05 & [+2.43, +3.68] & 14/0/86 \\
All & 17.828 & 17.389 & +2.60 & [+2.35, +2.84] & 70/23/407 \\
\bottomrule
\end{tabular}
\end{table}

Eight-view augmentation narrows the gap. With eight views, the mean gap on test500 is +1.13\% [+1.02, +1.25], E133 wins 66, ties 51 and loses 383 of 500 instances, and the ratio of mean costs is +1.05\% (17.400 against 17.219). By size, the gap is +0.46\%, +1.02\%, +1.32\%, +1.48\% and +1.39\% for $J = 20$, 40, 60, 80 and 100. Augmentation thus helps E133 more than G200, but G200 remains better at every size. On the validation set, G200 is 2.23\% better than E133 (17.309 against 17.703).

\subsection{Generalisation to larger instances}

Table~\ref{tab:gen} reports instances with 150 to 500 targets, which are beyond the training range. Under eight views, the mean gap to G200 is +0.72\%, and at $J = 300$ the interval includes zero. These gaps are smaller than the unaugmented test500 gaps, but the two protocols differ, so the tables are not directly comparable. Without augmentation, the mean gap on the same instances is +1.63\%. Augmentation helps both models, and it helps E133 more: the unaugmented mean is 36.801 for E133 and 36.336 for G200, and eight views reduce these to 35.827 and 35.644. Neither model was trained at these sizes. Because the two models differ in network and training length as well as in training signal, these results describe the two recipes; they do not show what the training signal alone contributes to transfer.

\begin{table}[t]
\centering
\caption{Generalisation set with eight views. Gap is E133 relative to G200.}
\label{tab:gen}
\small
\begin{tabular}{rrrrrr}
\toprule
$J$ & E133 & G200 & Gap (\%) & 95\% interval & Win/Tie/Loss \\
\midrule
150 & 26.003 & 25.795 & +0.96 & [+0.62, +1.31] & 23/0/77 \\
200 & 29.077 & 28.868 & +0.97 & [+0.56, +1.37] & 30/0/70 \\
300 & 36.945 & 36.875 & +0.38 & [$-$0.03, +0.78] & 42/0/58 \\
500 & 51.283 & 51.040 & +0.56 & [+0.23, +0.90] & 40/0/60 \\
All & 35.827 & 35.644 & +0.72 & [+0.53, +0.91] & 135/0/265 \\
\bottomrule
\end{tabular}
\end{table}

\subsection{Results on the literature benchmark}
\label{sec:real}

Table~\ref{tab:real} reports Real529 by size band. The mean gap to G200 is +0.21\%, with an interval that excludes zero. Gaps are below +0.53\% in every band up to $J = 100$. At $J = 200$, E133 is 1.25\% better than G200, with 21 wins out of 25. Ties are frequent at the smallest sizes, where 127 of 220 instances end at the same value, which suggests that both models often reach the same route. This holds for 25 instances from one source and lies outside the training range, so we treat it as an observation and not as evidence of general superiority: on the synthetic instances at the same size in Table~\ref{tab:gen}, E133 is behind.

\begin{table}[t]
\centering
\caption{Real529 with eight views, by size band. Gap is E133 relative to G200.}
\label{tab:real}
\small
\begin{tabular}{lrrrrrr}
\toprule
Band & $n$ & E133 & G200 & Gap (\%) & 95\% interval & Win/Tie/Loss \\
\midrule
$J = 10$--19 & 220 & 255.570 & 255.086 & +0.19 & [+0.12, +0.26] & 23/127/70 \\
$J = 20$--50 & 234 & 342.841 & 341.847 & +0.33 & [+0.21, +0.45] & 60/33/141 \\
$J = 100$ & 50 & 431.749 & 429.727 & +0.52 & [+0.13, +0.92] & 19/0/31 \\
$J = 200$ & 25 & 720.934 & 730.123 & $-$1.25 & [$-$1.78, $-$0.72] & 21/0/4 \\
All & 529 & 332.819 & 332.421 & +0.21 & [+0.13, +0.29] & 123/160/246 \\
\bottomrule
\end{tabular}
\end{table}

Table~\ref{tab:meta} compares both models with a reference value and with three metaheuristics that search over visiting orders and score every order with the same exact SOCP. The reference value of each instance is the best upper bound over all methods in \citet{li2025}. It is not a best known solution in general, which is why a few entries in the table are negative. \placeholder{Merge the best values of Vieira et al. (2026) into the reference; per-instance values are needed, since their preprint reports the Poikonen set as means by size.} Each metaheuristic was given one, three or ten times the single-core time that E133 needs on that instance. At equal time, iterated local search (ILS) and variable neighbourhood search (VNS) are worse than E133 overall (+2.93\%), but they are already better than both learned policies on the clustered instances of \citet{poikonen2020} (+1.30\% and +1.31\% against +2.10\% for E133 and +1.76\% for G200). At three times the time their mean gaps (+1.63\% and +1.65\%) are still larger than that of E133, although ILS is better than E133 on more instances than it is worse (251 against 180). At ten times the time they are better than both learned policies, at +0.73\% and +0.78\% against +1.35\% for E133 and +1.14\% for G200. Clustered instances differ from the uniform distribution used for training, which is consistent with the weakness of both learned policies there. Simulated annealing (SA) remains behind at all budgets, and its gap hardly changes between one and three times the time (+4.53\% and +4.52\%), which indicates that our SA implementation is poorly calibrated for this problem; we report its ten-times row for completeness and draw no conclusion from it. On average, E133 and G200 need 0.78\,s of wall time per instance, or 10.3\,s on a single core, on the same hardware. The time grows with size, from under 0.1\,s of wall time at $J = 10$ to about 5.6\,s at $J = 200$, with 16 parallel solver workers.

The subset averages in Table~\ref{tab:meta} hide a strong effect of size. Measured against the same reference, E133 and G200 are +4.50\% and +4.02\% above it at $J = 100$, and +5.91\% and +7.09\% at $J = 200$. The two-stage iterative heuristic (TSIH) of \citet{li2025} is +0.30\% above the reference over all 529 instances at a mean of 20.2\,s per instance, and +0.51\% and +0.46\% at $J = 100$ and $J = 200$, at 45.9\,s and 164\,s. These figures are taken from \citet{li2025}, whose hardware differs from ours, and the times are not normalised. The learned policies are therefore far from the strongest published heuristic at the largest sizes, which is also where their advantage in wall time is largest. \placeholder{After the Vieira merge, count instances on which a method in this paper is below the reference by more than rounding (reference values have two decimals), recheck each with an independent feasibility check, and report them. No new best solutions are claimed until then.}

\begin{table}[t]
\centering
\caption{Real529 gap to the reference value $R$ (best upper bound over all methods in \citet{li2025}), $100\,(\mathrm{UB}-R)/\mathrm{UB}$, where UB is the upper bound found by the method, by subset. Budgets are multiples of the single-core time of E133 on each instance. Our SA implementation is poorly calibrated (see text).}
\label{tab:meta}
\small
\resizebox{\textwidth}{!}{\begin{tabular}{lrrrrrr}
\toprule
Method & Gam72 & \begin{tabular}[c]{@{}c@{}}Erd132\\ 16--20\end{tabular} & \begin{tabular}[c]{@{}c@{}}Erd132\\ 25--50\end{tabular} & \begin{tabular}[c]{@{}c@{}}Poi325\\ Regular\end{tabular} & \begin{tabular}[c]{@{}c@{}}Poi325\\ Cluster\end{tabular} & All \\
\midrule
E133 & +0.19 & +0.24 & +0.89 & +1.75 & +2.10 & +1.35 \\
G200 & +0.01 & +0.08 & +0.80 & +1.59 & +1.76 & +1.14 \\
ILS, $1\times$ & +1.20 & +1.40 & +4.21 & +5.03 & +1.30 & +2.93 \\
VNS, $1\times$ & +1.20 & +1.40 & +4.21 & +5.04 & +1.31 & +2.93 \\
ILS, $3\times$ & +0.71 & +0.80 & +3.08 & +2.79 & +0.36 & +1.63 \\
VNS, $3\times$ & +0.69 & +0.76 & +3.14 & +2.84 & +0.36 & +1.65 \\
ILS, $10\times$ & +0.33 & +0.27 & +1.90 & +1.22 & $-$0.02 & +0.73 \\
VNS, $10\times$ & +0.20 & +0.34 & +1.89 & +1.40 & $-$0.03 & +0.78 \\
SA, $10\times$ & +0.53 & +1.27 & +5.15 & +5.16 & +1.84 & +3.14 \\
\bottomrule
\end{tabular}}
\end{table}

\subsection{Comparison with a general-purpose exact solver}

To set the scale of the problem, we also solved the full mixed-integer SOCP (Model+ of \citet{li2025}) with Gurobi, with a time limit of 3,600\,s per instance and one thread. \placeholder{Report status, upper bound, lower bound, gap and time by $J$ for Real529 and test500. Both runs have finished; results are in outputs/full\_gurobi\_20261030.}

\subsection{Training efficiency}
\label{sec:efficiency}

We compare training budgets in two ways. The first holds the solver budget fixed. POMO trained from scratch under the G200 recipe on our workstation was stopped once it had used the same 3.84 million exact solves as E133, which took it 264,512 instances. Its mean validation cost is 4.44\% higher than that of E133 (18.488 against 17.703), and E133 is better on 1,406 of the 2,000 validation instances. On Real529, with eight views, its mean per-instance gap to E133 is +3.79\% [+3.28, +4.30] and its mean cost is 5.67\% higher; E133 is better on 354 instances, tied on 110 and worse on 65. The gap grows with size, to 11.7\% at $J = 100$ and 24.0\% at $J = 200$. After 13.2 million solves the same run is still 2.2\% behind E133 on validation (18.089), so under the baseline recipe more than 3.4 times the solver budget of E133 is needed to match it. Figure~\ref{fig:efficiency} shows validation cost against cumulative solves and instances. The second comparison is with G200 at its full budget, which is better than E133 (Section~\ref{sec:synthetic}).

\begin{figure}[t]
\centering
\includegraphics[width=\textwidth]{data_efficiency.pdf}
\caption{Validation cost (mean over 2,000 instances, $J$ POMO starts, exact costs) against cumulative exact solves (left) and training instances (right). Proposed: the run that produced E133, every evaluated epoch. POMO with exact rewards: trained from scratch under the G200 recipe on our workstation. G200 is shown as a single point at its training budget.}
\label{fig:efficiency}
\end{figure}

Table~\ref{tab:budget} lists the full budgets. E133 used 15 times fewer training instances and 78 times fewer Gurobi calls in training than G200. The second ratio is the product of two factors: 15 times fewer instances and 5.2 times fewer calls per instance. G200 solved $J/4$ routes per training instance, so a baseline that trained with $J$ starts, as both models use at test time, would have needed about four times as many calls. These ratios stop at E133 and count training only. The run that produced E133 continued to epoch 200, using 6,239,185 training calls and 400,000 validation calls, and selecting E133 among its checkpoints by full POMO validation took 5,175,733 further solves. As run, the total is therefore 11,814,918 calls. The training log of G200 counts 338,545,920 calls when the epochs that were restarted after interruptions are counted each time they ran, about 29 times our total. G200 used its final epoch and needed no checkpoint selection, and its log records no validation calls. The test-time cost is unchanged in our measurements because it is dominated by exact solves of the POMO routes: on Real529, decoding takes 0.12\,s per instance on the GPU and Gurobi takes about 10.2\,s of single-core time for either model. CPU time for G200 is an estimate, obtained by multiplying its call count by a mean solve time of 0.031\,s measured at about $J = 60$. The solves of E133 took 0.081\,s on average over its mix of sizes, so at that rate the estimate for G200 would be about 6,800 core-hours; E133 also spent 67.6 core-hours on table search. Wall time for G200 comes from a different machine with 128 reward workers, so the ratio of wall times shows the order of magnitude only. Solver calls and instances are exact counts and form the primary comparison.

Table~\ref{tab:calls} breaks down the 3.84 million calls of E133. Anchor and target solves contribute one call per instance each. Expert verification, including near-miss checks, contributes 0.37 calls per instance. Refresh rounds add 0.52 calls per instance on average across all sizes, and more than one per instance at the end of the run, when more than half of the batches have $J \ge 60$ and refresh is active. A share of 23.7\% of refresh solves produced a better expert, so the refresh budget has a measurable return but also a large waste. Near-miss checks were attempted 364,974 times.

\begin{table}[t]
\centering
\caption{Exact solver calls of E133, epochs 1 to 133, by purpose.}
\label{tab:calls}
\small
\begin{tabular}{lrr}
\toprule
Purpose & Calls & Per instance \\
\midrule
Anchor solve & 1,330,000 & 1.00 \\
Target reward solve & 1,330,000 & 1.00 \\
Expert and near-miss verification & 493,128 & 0.37 \\
Refresh rounds & 686,842 & 0.52 \\
Total & 3,839,970 & 2.89 \\
\bottomrule
\end{tabular}
\end{table}

Table~\ref{tab:dynamics} shows how the loop behaves during training. In the first epoch, almost any search improves on an untrained policy. By epoch 40 the plain 2-opt search improves only 9\% of the anchors, by 0.38\% on average. The ramp then lifts the share of improved experts above one third, which is the intended effect of the progressive search. Over the same period, the lowest-bound policy route has a bound that falls from 92\% to 34\% above its exact cost, which is consistent with policy routes moving closer to the anchors as training proceeds. The bound remains loose at epoch 133, yet training continues to make progress; Section~\ref{sec:loose} discusses what this does and does not show.

\begin{table}[t]
\centering
\caption{Training dynamics of the proposed method. Improved share is the fraction of instances whose expert beats the anchor. Bound excess compares the lowest bound among policy routes with the exact cost of that route.}
\label{tab:dynamics}
\small
\begin{tabular}{rrrrr}
\toprule
Epoch & Anchor cost & Expert cost & Improved share (\%) & Bound excess (\%) \\
\midrule
20 & 19.986 & 19.779 & 12.1 & 92.4 \\
40 & 19.205 & 19.096 & 9.0 & 70.7 \\
60 & 19.146 & 18.798 & 24.4 & 57.5 \\
80 & 18.993 & 18.640 & 33.0 & 57.9 \\
100 & 18.850 & 18.603 & 34.9 & 44.2 \\
133 & 18.522 & 18.288 & 37.9 & 33.6 \\
\bottomrule
\end{tabular}
\end{table}

\begin{table}[t]
\centering
\caption{Training budgets of G200 and E133.}
\label{tab:budget}
\small
\resizebox{\textwidth}{!}{\begin{tabular}{lrrr}
\toprule
 & G200 & E133 & Ratio \\
\midrule
Training instances (to E133) & 20,000,000 & 1,330,000 & 15 \\
Gurobi calls in training (to E133) & 299,437,760 & 3,839,970 & 78 \\
Gurobi calls as run\textsuperscript{a} & 338,545,920 & 11,814,918 & 29 \\
Gurobi CPU (core-hours, to E133) & about 2,580 (estimated at 0.031\,s per call) & 86.7 (0.081\,s per call) & about 30 \\
Table search CPU (core-hours, to E133) & -- & 67.6 & -- \\
Wall time (hours) & about 215.5 (128 workers, other machine) & 19.2 to E133, 28.3 to E200 (32 workers) & about 11 \\
Test time per Real529 instance (s) & 0.78 & 0.78 & 1 \\
\bottomrule
\end{tabular}}
\par\smallskip\footnotesize\raggedright \textsuperscript{a}G200: all training epochs, with epochs restarted after interruptions counted each time they ran; it used its final epoch, so no checkpoint selection, and its log records no validation calls. E133: training to E200 (6,239,185), in-training validation (400,000) and checkpoint selection on validation2000 (5,175,733).
\end{table}

\placeholder{Add the 24-hour training curve: validation cost against wall time for POMO with exact rewards (Group A, running since 9 November 2026) and for the proposed method on the same hardware; confirm the plan for the proposed-method curve.}

\subsection{Ablation study}
\label{sec:ablation}

The ablations measure what individual components contribute at a fixed point in training; they report component effects and are not offered as a decomposition of the saving. Variants B and D train from scratch to one million instances (epoch 100), with the configuration and seed of the proposed method otherwise unchanged. They are compared with Group C, the proposed method at the same point (validation mean 17.772 at epoch 100), on test500 with eight views, by size. We fixed the decision rule before seeing the results: a component is said to have an effect only if the paired per-instance difference between C and the variant has a 95\% interval that excludes zero. Conclusions apply to epoch 100, before the reported checkpoint E133, and each variant is a single run.
\begin{itemize}
\item Group A retrains POMO with exact rewards from scratch under the G200 recipe, including its $J/4$ training starts, on the same hardware, for 24 hours. Its checkpoints at the instance and solver budgets of E133 extend the equal-budget comparison of Section~\ref{sec:efficiency}.
\item Group B lets the expert search evaluate candidate moves exactly instead of on the table. Because exact evaluation is expensive, the search is also reduced: it starts from the anchor only, runs first-improvement 2-opt with at most 100 exact solves per instance, and uses no LKH and no refresh. The reward for policy routes still uses the anchor table. Group B therefore measures the role of the table inside the expert search, together with the effect of a smaller search, and not its role as a reward. To epoch 100 it made 101.8 million exact solves, against 2.3 million for C, and every instance used its full search budget; its training took 4.4 times as long (48.6 against 11.0 hours).
\item Group D reduces the reference pool to the greedy route of the policy from the first start, which receives the local-search budget of the 30 starts. It varies the diversity of the starts of the expert search, together with their quality. It does not vary the diversity of the POMO routes discussed in Section~\ref{sec:limits}.
\end{itemize}
\placeholder{Report quality by size for B, C and D on test500 with eight views (queued after Group A), with the decision rule above, and Group A's checkpoints.}

\section{Discussion and implications for practice}
\label{sec:discussion}

\subsection{When the approach is worthwhile}

A trained policy moves most of the computing effort into training, but not all of it: at test time E133 still solves one SOCP per distinct decoded route, about 10.3\,s of single-core time per Real529 instance, or 0.78\,s of wall time with 16 solver workers. Within the range of speeds and endurance seen in training, the policy can be reused without retraining, because both enter as input features. The same holds for G200, which gives better plans at the same test time. The smaller training budget therefore matters when a learned solver has to be trained from scratch, for example for a variant with additional constraints or a different spatial distribution of targets. In that situation the solver budget of training decides whether a learned solver is affordable at all, and our results show that a useful policy can be obtained with a few million solves instead of hundreds of millions. We did not test training on a changed problem.

Table~\ref{tab:breakeven} sets the computing cost of the options against their plan quality on Real529. The training CPU of E133 to epoch 133, 154.3 core-hours of Gurobi and table search, equals the additional single-core time of ILS at ten times the budget on about 6,000 instances, at about 93\,s per instance. Below that number of plans, ILS with ten times the budget uses less CPU in total and gives better plans. Above it, the policy uses less CPU, but its plans are about 0.6 percentage points further from the reference on average, and much further at $J \ge 100$ (Section~\ref{sec:real}). Where G200 is already available it is preferable to E133, since its test cost is the same and its plans are better. The comparison treats training CPU and planning CPU as equal in value and ignores GPU time.

\begin{table}[t]
\centering
\caption{Computing cost and plan quality on Real529. Single-core time per instance for the metaheuristics is the stated multiple of the mean single-core time of E133. Gap is to the reference value of Table~\ref{tab:meta}.}
\label{tab:breakeven}
\small
\resizebox{\textwidth}{!}{\begin{tabular}{lrrrr}
\toprule
Option & Training solver calls & Training CPU (core-hours) & Single-core time per instance (s) & Gap (\%) \\
\midrule
G200 & 299,437,760 & about 2,580 (estimated) & about 10.3 & +1.14 \\
E133 & 3,839,970 (about 11.8 million as run) & 154.3 & about 10.3 & +1.35 \\
ILS, $3\times$ & -- & -- & about 31 & +1.63 \\
ILS, $10\times$ & -- & -- & about 103 & +0.73 \\
\bottomrule
\end{tabular}}
\end{table}

At the time budgets we tested, the advantage of the policy over a metaheuristic is limited. ILS with the same single-core time is behind E133 overall but ahead on clustered instances, and with three times the time it is better on more instances than it is worse (251 against 180), although its mean gap is larger. The clear advantage of the policy lies in wall time when solver workers can run in parallel, which we did not give the metaheuristics. A route from the policy could also start a metaheuristic, and the table could be used for search at test time; we did not test either.

\subsection{What the loose bound does and does not show}
\label{sec:loose}

The bound is far from tight on routes that differ from the anchor. The lowest bound among policy routes exceeds the exact cost of that route by 33.6\% at epoch 133 (Table~\ref{tab:dynamics}), and near-miss routes, which come out of local search, carry bounds 14.6\% above the anchor cost on average although 40.8\% of them are exactly better (Section~\ref{sec:search}). The bound is thus not accurate near the anchor at the resolution that matters for acceptance, and most late-stage experts are accepted through near-miss checks and refresh rounds. In the recipe it serves mainly as a reward for routes that are never solved and as a ranking that decides which few routes receive an exact solve, while the exactness of every label keeps errors in the bound out of imitation learning.

Two experiments bear on whether a tighter or less biased signal would help, and neither settles the question. In a 30-epoch continuation from E133, taking the minimum over all tables available for each route used about two additional tables per instance on the affected batches, tightened the bound by only 0.32\% and left validation quality unchanged. This is a small change against an excess of about one third, so it does not show that tightness has low value. In an earlier 20-epoch study from scratch, rewards from exact solves of a random subset of 4 to 8 routes per instance point in a direction much closer to the all-exact advantage (an expected cosine of about 0.85 to 0.95, against about 0.6 for the hybrid reward). They cost three to four times the Gurobi CPU and did not improve quality by epoch 20: they were worse in one pair of runs and level in the other. That study is short and single-seed. The hybrid reward is also biased, not only noisy (Section~\ref{sec:learning}), and whether a tighter or unbiased signal would close the remaining gap is open. Two refinements that keep the guarantee are available and untested: letting each target take its sortie from a different solved solution, chosen by a dynamic programme along the route, and letting each target exchange its take-off and landing points, chosen by a two-state dynamic programme along the route.

\subsection{Limitations}
\label{sec:limits}

Six limitations apply. First, E133 remains behind G200 on synthetic instances from the training distribution, by +1.13\% with eight views and by up to about 3.5\% at some sizes without augmentation. We have not identified the cause, and several explanations remain open. (i) Route diversity: in a diagnostic on a later checkpoint of the same run (epoch 163), the edge overlap between POMO routes was higher than for G200 at every size tested (0.64, 0.67 and 0.62 against 0.48, 0.55 and 0.53 at $J = 20$, 60 and 100), and augmentation, which adds diversity at test time, helps E133 more than G200. Against this explanation, the overlap difference shrinks with size while the gap grows, and an earlier diagnostic found G200 no more diverse than a predecessor of the method that had a similar quality deficit, with the difference lying in the quality of individual routes. (ii) Pull towards the anchor: the bound is tightest at the anchor and the hybrid reward favours the solved route (Section~\ref{sec:learning}), so the reward may draw the policy towards anchor-like routes. In matched 20-epoch runs from the same initialisation, replacing the bound-based reward with random-subset exact rewards lowered the overlap at $J = 20$ from 0.374 to 0.330 and from 0.379 to 0.307. (iii) Training length and recipe: E133 saw 15 times fewer instances and uses a different network, and G200 trained with $J/4$ starts. Continuations from E133 that changed the bound or the imitation weight moved neither overlap nor quality appreciably, so they do not discriminate between (i) and (ii). Separating these explanations requires controlled runs that we have not performed. Second, every model is one training run with one seed, and the reported intervals cover variation over test instances only, not over training runs. Large differences, such as the 4.4\% gap at equal solver budget, are unlikely to reverse with another seed, but small ones, such as the +0.21\% gap on Real529, should be read as closeness and not as a ranking. Validation means across checkpoints 120 to 150 of the E133 run varied by 0.18\%; this is fluctuation between checkpoints, not variance between seeds. Third, Real529 is not held out and covers only two parameter settings, both at the edges of the training range. Fourth, the comparison on time depends on hardware and implementation, which is why we treat solver calls as the primary measure. Fifth, the settings of the progressive search were set by development runs and not by a formal search, so the reported performance may not be the best attainable for the method. Sixth, the bound depends on the separable structure of the CVTSP. It extends to a variant only if sortie constraints remain independent of the order, which should hold for several sorties of fixed composition but fails if, for instance, target time windows couple the sorties.

\section{Conclusion}
\label{sec:conclusion}

This paper asked whether a DRL policy for the CVTSP needs an exact solver reward for every route. One optimal fixed-order solution yields an upper bound on all visiting orders, computed by summing table entries, and a training recipe built on this bound, with 2.89 exact solves per training instance and exact verification of every imitation label, obtains a useful policy at a small fraction of the usual solver budget. With the same 3.84 million solves, POMO trained on exact rewards is on average 3.8\% worse per instance than the policy from our recipe on the literature benchmark, and on validation it needs more than 3.4 times as many solves to match it. The fully trained baseline used 78 times more solves in training, or about 29 times more when both runs are counted as run, and it remains better: by +0.21\% on the literature benchmark and +0.72\% on larger instances, where the two are close, and by +1.13\% with augmentation on synthetic instances from the training distribution. Iterated local search and variable neighbourhood search given ten times the single-core time of the policy, and the strongest published heuristic, obtain better solutions, most clearly at 100 and 200 targets. The recipe is therefore a way to obtain a learned CVTSP solver with a small solver budget, not a better solver.

Several directions follow. The recipe should be tested where training from scratch is needed, such as variants with further constraints or other target distributions. The training signal can be tightened while keeping its guarantee, through the per-target combination of tables and the reversal symmetry. The explanations for the remaining gap, route diversity, the pull of a biased reward towards the anchor and the difference in training length, need to be separated by controlled runs. The table could also support search at test time or warm starts for metaheuristics. The bound requires only that the sorties of one solution remain feasible for other orders, but whether it is useful elsewhere also depends on how loose it is, which we could not predict in advance. Several carriers \citep{fahradyan2020} and several UAVs on one vessel \citep{ren2026} raise different questions, since sorties may then overlap in time and the leg structure changes.

\section*{CRediT authorship contribution statement}
\placeholder{To be completed by the authors.}

\section*{Declaration of competing interest}
\placeholder{To be completed by the authors.}

\section*{Funding}
\placeholder{To be completed by the authors.}

\section*{Data availability}
\placeholder{State where the code, generated instances and per-instance results are available. Real529 instances are from the cited sources.}

\section*{Ethics statement}
This study used computational experiments on synthetic and published benchmark instances and involved no human participants or personal data.

\section*{Declaration of generative AI use}
\placeholder{The authors should confirm this statement.} During the preparation of this work the authors used Claude (Anthropic) to assist with literature search and manuscript drafting. The authors reviewed and edited the content and take full responsibility for the publication.

\bibliographystyle{elsarticle-harv}
\bibliography{refs}

\end{document}
